\documentclass[a4paper,11pt]{scrartcl}
\usepackage[margin=2.35cm]{geometry}
\usepackage{fontspec,unicode-math}
\setmainfont{Libertinus Serif}\setsansfont{Libertinus Sans}\setmonofont{Libertinus Mono}\setmathfont{Libertinus Math}
\usepackage{microtype,amsmath,booktabs,longtable,array,graphicx,needspace}
\usepackage[hidelinks]{hyperref}
\setkomafont{disposition}{\normalfont\bfseries}
\setlength{\emergencystretch}{2em}
\renewcommand{\thesection}{S\arabic{section}}
\renewcommand{\thetable}{S\arabic{table}}
\title{Supplementary material\\\large Sequence and family context shape ETV4 representations in protein language models}
\author{}\date{}
\begin{document}\maketitle
This supplement documents the existing computational comparisons, controls and provenance. It does not add model experiments. Tables retain all tested blocks and unfavourable outcomes; the accompanying TSV files retain full numerical precision.
\section{Biological regions, constructs and observations}
Human ETV4 uses native UniProt P43268 sequence version 3. Region names follow the manuscript. NID is predominantly disordered; CID includes the inhibitory H4 helix and is not wholly disordered. ETS + CID (337--436) also includes the two-residue junction. Shortened CID denotes 421--430. The published observations are from Currie et al., \emph{Nucleic Acids Research} 45, 2223--2241 (2017), \href{https://doi.org/10.1093/nar/gkx068}{doi:10.1093/nar/gkx068}, including Supplementary Tables S1, S3 and S10. They are biological reference observations, not outputs of the protein models.

\Needspace{13\baselineskip}
\begingroup\small
\begin{longtable}{@{}ll@{}}
\caption{Region definitions used throughout the study.}\\
\toprule Region & Native residues\\\midrule\endfirsthead
\toprule Region & Native residues\\\midrule\endhead
N-terminal extension & 1--164\\
NID & 165--336\\
Junction & 337--338\\
Folded ETS domain & 339--420\\
CID & 421--436\\
C-terminal tail & 437--484\\
\bottomrule\end{longtable}
\endgroup
Seven WT constructs and five mutant genotypes in each of two contexts give 17 distinct native input sequences. Y401N/Y403G and E404K affect ETS; E423K, E425K and F428A/F432A affect CID. The model-input registry includes the entire input sequence even when a substituted residue is outside the readout mask. Y401N/Y403G weakens DNA binding in ETS + CID and relieves NID-associated inhibition in the published comparison. E404K increases affinity while inhibition is reported as preserved. The CID substitutions are controls with binding and inhibition reported as preserved; some source statements lack displayed quantitative affinity data. We do not infer numerical equivalence from those statements.

Full length versus ETS + CID changes the N-terminal extension, NID and the C-terminal tail. It does not isolate NID. The ProFam WT comparison instead adds NID to ETS + CID, keeping the same prompt before each target.

\Needspace{14\baselineskip}
\begingroup\small
\begin{longtable}{@{}lll@{}}
\caption{Published WT affinities: means and reported SEM, converted to nM. Lower $K_D$ means tighter DNA binding. Replicate counts remain in the source TSV.}\\
\toprule Construct & $K_D$ (nM) & SEM (nM)\\\midrule\endfirsthead
\toprule Construct & $K_D$ (nM) & SEM (nM)\\\midrule\endhead
ETS + shortened CID & 0.061 & 0.006\\
ETS + CID & 0.120 & 0.020\\
ETS + CID + tail & 0.110 & 0.020\\
NID + ETS + shortened CID & 0.120 & 0.010\\
NID + ETS + CID & 0.600 & 0.100\\
NID + ETS + CID + tail & 0.660 & 0.090\\
Full length & 0.830 & 0.080\\
\bottomrule\end{longtable}
\endgroup
\section{Exact readout masks and model measurements}
The structural ETS domain is 339--420 (82 residues). The shared readout window 337--430 has 94 residues: two junction, 82 ETS and ten CID positions. The ESM mutation readout excludes the union of substituted positions inside this window: 401, 403, 404, 423, 425 and 428. Its 88 positions comprise two junction, 79 ETS and seven CID positions. F432 remains in the input but is outside the window. The WT factorial and ProFam readouts use all 94 positions. The complete masks are also machine-readable in the original run plan.

\noindent\textbf{ESM mutation mask (88 positions):} 
337, 338, 339, 340, 341, 342, 343, 344, 345, 346, 347, 348, 349, 350, 351, 352, 353, 354, 355, 356, 357, 358, 359, 360, 361, 362, 363, 364, 365, 366, 367, 368, 369, 370, 371, 372, 373, 374, 375, 376, 377, 378, 379, 380, 381, 382, 383, 384, 385, 386, 387, 388, 389, 390, 391, 392, 393, 394, 395, 396, 397, 398, 399, 400, 402, 405, 406, 407, 408, 409, 410, 411, 412, 413, 414, 415, 416, 417, 418, 419, 420, 421, 422, 424, 426, 427, 429, 430.\par
At each position, the post-block vector is divided by its Euclidean norm. For a vector difference $w$, RMS denotes $\sqrt{|S|^{-1}\sum_{i\in S}\|w_i\|_2^2}$. Let $u$ be mutant minus WT in ETS + CID and $v$ mutant minus WT in full length. Then $D=\mathrm{RMS}(u)$ and $M=\mathrm{RMS}(v-u)$. $M$ is not a difference of scalar magnitudes. $D$ is not a binding energy. $M/D$ is a prespecified diagnostic without an established biological normalization.

The WT contrast is $J=z(NC)-z(N)-z(C)+z(0)$, with $0$ = ETS + shortened CID, $N$ = NID + ETS + shortened CID, $C$ = ETS + CID, and $NC$ = NID + ETS + CID. ProFam measures $q_P=z(NC\mid P)-z(C\mid P)$ for prompt $P$. Its left-to-right architecture prevents ETS residues from seeing downstream CID and the whole readout from seeing 431--436. The resulting zero ProFam $J$ is an architectural check.
\section{Numerical allowances and every fixed ESM block}
Numerical qualification used artificial inputs and froze $\eta$ before biological inference. Denominators must exceed $4\eta$. Ratio bounds are $L=\max(0,M-8\eta)/(D+4\eta)$ and $U=(M+8\eta)/(D-4\eta)$. They are sensitivity bounds under the measured implementation discrepancy, not confidence intervals or a guarantee covering all numerical error.

\Needspace{16\baselineskip}
\begingroup\small
\begin{longtable}{@{}lllll@{}}
\caption{Frozen numerical allowances, all displayed in units of $10^{-3}$. Exact values remain in the execution configurations.}\\
\toprule Model & Block & $\eta$ & $4\eta$ & $8\eta$\\\midrule\endfirsthead
\toprule Model & Block & $\eta$ & $4\eta$ & $8\eta$\\\midrule\endhead
650M & 17 & 0.447 & 1.790 & 3.579\\
650M & 22 & 0.364 & 1.458 & 2.916\\
650M & 33 & 0.190 & 0.761 & 1.522\\
3B & 19 & 0.122 & 0.489 & 0.979\\
3B & 24 & 0.171 & 0.682 & 1.365\\
3B & 36 & 0.118 & 0.470 & 0.941\\
ProFam & 8 & 0.067 & 0.270 & 0.540\\
ProFam & 11 & 0.055 & 0.221 & 0.442\\
ProFam & 16 & 0.004 & 0.018 & 0.036\\
\bottomrule\end{longtable}
\endgroup

\begingroup\small
\begin{longtable}{@{}lllllll@{}}
\caption{All fixed-block ESM mutation measurements. Bounds use the model/block-specific allowance.}\\
\toprule Model/block & Genotype & $M$ & $D$ & $M/D$ & $L$ & $U$\\\midrule\endfirsthead
\toprule Model/block & Genotype & $M$ & $D$ & $M/D$ & $L$ & $U$\\\midrule\endhead
650M/17 & Y401N/Y403G & 0.024 & 0.054 & 0.441 & 0.362 & 0.525\\
650M/17 & E404K & 0.011 & 0.019 & 0.613 & 0.385 & 0.890\\
650M/17 & E423K & 0.009 & 0.020 & 0.446 & 0.244 & 0.688\\
650M/17 & E425K & 0.011 & 0.019 & 0.551 & 0.335 & 0.810\\
650M/17 & F428A/F432A & 0.019 & 0.038 & 0.501 & 0.389 & 0.626\\
650M/22 & Y401N/Y403G & 0.024 & 0.054 & 0.433 & 0.370 & 0.500\\
650M/22 & E404K & 0.012 & 0.018 & 0.652 & 0.452 & 0.887\\
650M/22 & E423K & 0.009 & 0.019 & 0.500 & 0.320 & 0.711\\
650M/22 & E425K & 0.011 & 0.019 & 0.602 & 0.416 & 0.819\\
650M/22 & F428A/F432A & 0.019 & 0.038 & 0.505 & 0.413 & 0.605\\
650M/33 & Y401N/Y403G & 0.062 & 0.125 & 0.497 & 0.482 & 0.512\\
650M/33 & E404K & 0.029 & 0.058 & 0.501 & 0.469 & 0.534\\
650M/33 & E423K & 0.026 & 0.046 & 0.568 & 0.526 & 0.611\\
650M/33 & E425K & 0.027 & 0.044 & 0.626 & 0.581 & 0.673\\
650M/33 & F428A/F432A & 0.048 & 0.078 & 0.618 & 0.592 & 0.643\\
3B/19 & Y401N/Y403G & 0.026 & 0.057 & 0.447 & 0.427 & 0.468\\
3B/19 & E404K & 0.010 & 0.022 & 0.476 & 0.421 & 0.533\\
3B/19 & E423K & 0.008 & 0.018 & 0.421 & 0.357 & 0.489\\
3B/19 & E425K & 0.008 & 0.018 & 0.451 & 0.385 & 0.519\\
3B/19 & F428A/F432A & 0.017 & 0.031 & 0.554 & 0.515 & 0.594\\
3B/24 & Y401N/Y403G & 0.033 & 0.075 & 0.436 & 0.414 & 0.458\\
3B/24 & E404K & 0.012 & 0.025 & 0.478 & 0.412 & 0.548\\
3B/24 & E423K & 0.009 & 0.021 & 0.438 & 0.360 & 0.521\\
3B/24 & E425K & 0.009 & 0.021 & 0.456 & 0.377 & 0.540\\
3B/24 & F428A/F432A & 0.020 & 0.035 & 0.559 & 0.511 & 0.609\\
3B/36 & Y401N/Y403G & 0.084 & 0.147 & 0.571 & 0.563 & 0.580\\
3B/36 & E404K & 0.031 & 0.058 & 0.544 & 0.523 & 0.565\\
3B/36 & E423K & 0.024 & 0.042 & 0.572 & 0.544 & 0.601\\
3B/36 & E425K & 0.024 & 0.040 & 0.612 & 0.582 & 0.643\\
3B/36 & F428A/F432A & 0.047 & 0.065 & 0.718 & 0.698 & 0.737\\
\bottomrule\end{longtable}
\endgroup
The original directional rule required resolved denominators, raw-$M$ separation, ratio-bound separation and native-order separation for Y401N/Y403G. The overall specificity classification remained unmet in both models at every fixed block. The 3B native-order condition passed without rescuing ratio-bound separation. These failures do not establish absence of regulatory information.
\section{Sequence-order controls and seeds}
Twenty full-length controls shuffle residues 1--336 (N-terminal extension + NID), preserving composition, length and endpoints, with each shuffled prefix shared by WT and every mutant. Seeds are 13001--13020. WT factorial controls shuffle NID with seeds 11001--11020 and the six-residue CID completion with seeds 12001--12020, preserving endpoints. Only four CID-completion positions are permutable; repeated orders and one native-order draw are retained. Exceedance counts are descriptive, not $p$-values.

\Needspace{13\baselineskip}
\begingroup\small
\begin{longtable}{@{}llll@{}}
\caption{Y401N/Y403G fixed native-order test. Margin and required allowance are in units of $10^{-3}$. The unrounded comparison determines the outcome.}\\
\toprule Model/block & Margin & Required $8\eta$ & Pass\\\midrule\endfirsthead
\toprule Model/block & Margin & Required $8\eta$ & Pass\\\midrule\endhead
650M/17 & -0.620 & 3.579 & False\\
650M/22 & -0.875 & 2.916 & False\\
650M/33 & 1.521 & 1.522 & False\\
3B/19 & 1.698 & 0.979 & True\\
3B/24 & 1.916 & 1.365 & True\\
3B/36 & 6.528 & 0.941 & True\\
\bottomrule\end{longtable}
\endgroup
At the primary block of each ESM model, ten of twenty WT factorial controls equal or exceed the native RMS($J$). This is a different contrast from mutation-dependent native-order sensitivity. All individual WT and mutation-control values are retained in the data files.
\section{Protein prompts and all ProFam response magnitudes}
Five pairs selected from five reviewed ETV1/ETV5 entries (seed 20261110) were used in both orders. Controls are exactly length-matched nonhomologous prompts and composition-shuffled homologous prompts; the empty prompt supplies a baseline. Prompt-shuffle seeds are $14000+10s+m$, with set $s=1,\ldots,5$ and member $m=0,1$. There are 31 named settings. The ten nonhomologous settings contain only four distinct ordered prompts. They are not independent biological replicates. The exact sequences, ordering, hashes and token-prefix lengths are in \texttt{prompt-panel.json}.


\begingroup\small
\begin{longtable}{@{}llll@{}}
\caption{Prompt identity and lengths. Accessions for shuffled homologues identify the source sequences; sequences were shuffled before use. Prefix tokens include the model document/separator tokens.}\\
\toprule Setting & Accessions & AA lengths & Prefix tokens\\\midrule\endfirsthead
\toprule Setting & Accessions & AA lengths & Prefix tokens\\\midrule\endhead
empty & none & 0 & 2\\
hom\_1\_ab & P41164, P50549 & 477/477 & 958\\
hom\_1\_ba & P50549, P41164 & 477/477 & 958\\
nonhom\_1\_ab & P56270, Q969S3 & 477/477 & 958\\
nonhom\_1\_ba & Q969S3, P56270 & 477/477 & 958\\
homshuffle\_1\_ab & P41164, P50549 & 477/477 & 958\\
homshuffle\_1\_ba & P50549, P41164 & 477/477 & 958\\
hom\_2\_ab & P41161, Q2KIC2 & 510/477 & 991\\
hom\_2\_ba & Q2KIC2, P41161 & 477/510 & 991\\
nonhom\_2\_ab & Q8TBC5, Q969S3 & 510/477 & 991\\
nonhom\_2\_ba & Q969S3, Q8TBC5 & 477/510 & 991\\
homshuffle\_2\_ab & P41161, Q2KIC2 & 510/477 & 991\\
homshuffle\_2\_ba & Q2KIC2, P41161 & 477/510 & 991\\
hom\_3\_ab & P50549, Q2KIC2 & 477/477 & 958\\
hom\_3\_ba & Q2KIC2, P50549 & 477/477 & 958\\
nonhom\_3\_ab & P56270, Q969S3 & 477/477 & 958\\
nonhom\_3\_ba & Q969S3, P56270 & 477/477 & 958\\
homshuffle\_3\_ab & P50549, Q2KIC2 & 477/477 & 958\\
homshuffle\_3\_ba & Q2KIC2, P50549 & 477/477 & 958\\
hom\_4\_ab & P41164, Q2KIC2 & 477/477 & 958\\
hom\_4\_ba & Q2KIC2, P41164 & 477/477 & 958\\
nonhom\_4\_ab & P56270, Q969S3 & 477/477 & 958\\
nonhom\_4\_ba & Q969S3, P56270 & 477/477 & 958\\
homshuffle\_4\_ab & P41164, Q2KIC2 & 477/477 & 958\\
homshuffle\_4\_ba & Q2KIC2, P41164 & 477/477 & 958\\
hom\_5\_ab & P41164, Q9CXC9 & 477/510 & 991\\
hom\_5\_ba & Q9CXC9, P41164 & 510/477 & 991\\
nonhom\_5\_ab & Q969S3, Q8TBC5 & 477/510 & 991\\
nonhom\_5\_ba & Q8TBC5, Q969S3 & 510/477 & 991\\
homshuffle\_5\_ab & P41164, Q9CXC9 & 477/510 & 991\\
homshuffle\_5\_ba & Q9CXC9, P41164 & 510/477 & 991\\
\bottomrule\end{longtable}
\endgroup

\begingroup\small
\begin{longtable}{@{}llll@{}}
\caption{RMS response to adding NID under every prompt at all fixed ProFam blocks.}\\
\toprule Setting & Block 8 & Block 11 & Block 16\\\midrule\endfirsthead
\toprule Setting & Block 8 & Block 11 & Block 16\\\midrule\endhead
empty & 0.455 & 0.492 & 0.120\\
hom\_1\_ab & 0.409 & 0.411 & 0.144\\
hom\_1\_ba & 0.408 & 0.410 & 0.144\\
nonhom\_1\_ab & 0.446 & 0.506 & 0.124\\
nonhom\_1\_ba & 0.451 & 0.507 & 0.122\\
homshuffle\_1\_ab & 0.412 & 0.468 & 0.110\\
homshuffle\_1\_ba & 0.421 & 0.464 & 0.111\\
hom\_2\_ab & 0.404 & 0.397 & 0.119\\
hom\_2\_ba & 0.397 & 0.391 & 0.119\\
nonhom\_2\_ab & 0.445 & 0.497 & 0.125\\
nonhom\_2\_ba & 0.435 & 0.476 & 0.113\\
homshuffle\_2\_ab & 0.412 & 0.469 & 0.113\\
homshuffle\_2\_ba & 0.407 & 0.463 & 0.114\\
hom\_3\_ab & 0.408 & 0.410 & 0.146\\
hom\_3\_ba & 0.408 & 0.410 & 0.145\\
nonhom\_3\_ab & 0.446 & 0.506 & 0.124\\
nonhom\_3\_ba & 0.451 & 0.507 & 0.122\\
homshuffle\_3\_ab & 0.392 & 0.436 & 0.102\\
homshuffle\_3\_ba & 0.412 & 0.457 & 0.108\\
hom\_4\_ab & 0.410 & 0.411 & 0.144\\
hom\_4\_ba & 0.409 & 0.411 & 0.145\\
nonhom\_4\_ab & 0.446 & 0.506 & 0.124\\
nonhom\_4\_ba & 0.451 & 0.507 & 0.122\\
homshuffle\_4\_ab & 0.400 & 0.442 & 0.098\\
homshuffle\_4\_ba & 0.425 & 0.464 & 0.111\\
hom\_5\_ab & 0.398 & 0.390 & 0.118\\
hom\_5\_ba & 0.406 & 0.399 & 0.120\\
nonhom\_5\_ab & 0.435 & 0.476 & 0.113\\
nonhom\_5\_ba & 0.445 & 0.497 & 0.125\\
homshuffle\_5\_ab & 0.397 & 0.423 & 0.098\\
homshuffle\_5\_ba & 0.392 & 0.422 & 0.105\\
\bottomrule\end{longtable}
\endgroup
Homologous responses are smaller than matched nonhomologous responses in all ten settings at blocks 8 and 11. At block 16 they are larger in eight of ten settings. These counts and paired differences of magnitudes were prepared descriptively for the manuscript from existing outputs. They are not the same as vector differences between responses. Prompt-order and NID-shuffle values are included in the accompanying TSVs. This ordering reversal does not establish a reversal of biological inhibition.
\section{Execution and reproducibility records}
Original ESM-2 650M and ProFam inference used float32 on Apple MPS with CPU qualification; their primary blocks were 22 and 11. The original environment recorded Python 3.11.15, PyTorch 2.15.0, fair-esm 2.0.0, ProFam 0.1.8, Transformers 4.57.6 and NumPy 2.4.6. The 3B follow-up used an A100-SXM4-40GB, PyTorch 2.15.0+cu132, CUDA 13.2, TF32 disabled, batch size one and primary block 24. Model weights were fixed; summaries used float64. Inference seed was 19001.

There were 217 biological/control forwards for each ESM model and 277 for ProFam. The 3B artificial qualification was a separate 20-forward job. These counts are computational workloads, not independent samples. The follow-up was motivated by the 650M result and is not an independent confirmation or an isolated parameter-count comparison.

Separate arithmetic implementations previously checked 78,789 original and 28,311 follow-up values. Those checks were same-agent verification. Their receipts are included; packaging did not repeat model inference. The default repository workflow reconstructs figures and tables from saved summaries. Historical adapters, inputs, masks, seeds, configuration and checkpoint identities are retained separately. Raw activations are archived in Zenodo; model weights are not redistributed. A portable NumPy-only analysis recomputes the canonical summaries from saved activations without loading model weights. It verified all saved fields, including 107,102 numerical values, against the frozen summaries.
\section{Machine-readable data and release status}
The eleven TSV files cover all fixed ESM metrics, native-order margins, primary shuffled responses, WT factorial controls, all ProFam prompt responses, NID shuffles, vector contrasts, prompt-order contrasts, matched magnitude differences and published WT affinities. \texttt{DATA-DICTIONARY.md} defines their columns. Frozen source hashes and a reconstruction report accompany the repository.

Code, processed data and analysis scripts are publicly available in the \href{https://github.com/amoyag/etv4-context-representations}{ETV4 context representations repository}. The activation data, numerical qualification data and a versioned code snapshot are archived in \href{https://doi.org/10.5281/zenodo.24220637}{Zenodo (doi: 10.5281/zenodo.24220637)}. Original code is licensed under MIT; original derived data and written materials are licensed under CC BY 4.0. Third-party materials retain their own terms.
\end{document}
